\exercisesfor{4}

第 4 节的约定除非另有说明均有效。

\begin{enumerate}
\def\labelenumi{\arabic{enumi}.}
\item
  令\(\mathfrak{g}\)为幂零李代数，\(p\)（分别地，\(q\)）是满足\(\mathcal{C}^p\mathfrak{g} = \{0\}\)（分别地\(\mathcal{C}_q\mathfrak{g} = \mathfrak{g}\)）的最小整数。证明\(p = q + 1\)，且\(\mathcal{C}_i\mathfrak{g} \supset \mathcal{C}^{p - i}\mathfrak{g}\)。（使用命题 1 的论证。）
\item
  令\(\mathfrak{g}\)为维数为 1 的代数\(\mathfrak{h}\)与交换理想\(\mathfrak{g}'\)的半直积。令\(x \in \mathfrak{h}, x \neq 0\)，且\(u\)为\(\operatorname{ad}_{\mathfrak{g}} x\)在\(\mathfrak{g}'\)上的限制。
\end{enumerate}

\begin{enumerate}
\def\labelenumi{(\alph{enumi})}
\item
  g 为幂零的充要条件是 u 为幂零。
\item
  要使\(\mathfrak{g}\)的 Killing 形式为零，充要条件是\(\operatorname{Tr}(u^2) = 0\)。
\item
  由 (a) 和 (b) 推出，存在 Killing 形式为零的非幂零李代数。
\item
  由 (a) 推出，在满足\(\mathcal{C}^{p-1}\mathfrak{g} \neq \{0\}\)且\(\mathcal{C}^p\mathfrak{g} = \{0\}\)的幂零李代数中，可能有\(\mathcal{C}_i\mathfrak{g} \neq \mathcal{C}^{p-i}\mathfrak{g}\)。
\end{enumerate}

\begin{enumerate}
\def\labelenumi{\arabic{enumi}.}
\setcounter{enumi}{2}
\tightlist
\item
  \begin{enumerate}
  \def\labelenumii{(\alph{enumii})}
  \tightlist
  \item
    令\(\mathfrak{g}\)为幂零李代数，\(\mathfrak{z}\)为其中心，\(\mathfrak{b}\)为\(\mathfrak{g}\)的非零理想。证明\(\mathfrak{z} \cap \mathfrak{b} \neq \{0\}\)。（将\(\mathfrak{b}\)通过伴随表示视为\(\mathfrak{g}\)-模。）
  \end{enumerate}
\end{enumerate}

\begin{enumerate}
\def\labelenumi{(\alph{enumi})}
\setcounter{enumi}{1}
\tightlist
\item
  若在李代数\(\mathfrak{g}\)中，理想\(\mathfrak{b}\)包含于\(\mathcal{C}_{i+1}\mathfrak{g}\)但不包含于\(\mathcal{C}_i\mathfrak{g}\)，证明理想\(\mathfrak{b} \cap \mathcal{C}_k\mathfrak{g}\)在\(0 \leqslant k \leqslant i + 1\)所给范围内两两不同。（应用 (a)。）
\end{enumerate}

\begin{enumerate}
\def\labelenumi{\arabic{enumi}.}
\setcounter{enumi}{3}
\tightlist
\item
  令 g 为幂零李代数。
\end{enumerate}

\begin{enumerate}
\def\labelenumi{(\alph{enumi})}
\item
  \(\mathfrak{g}\)的每个子代数都是次不变的。（使用命题 3。）
\item
  令\(\mathfrak{b}\)为\(\mathfrak{g}\)的向量子空间，满足\(\mathfrak{b} + \mathcal{D}\mathfrak{g} = \mathfrak{g}\)。证明\(\mathfrak{g}\)的子代数由\(\mathfrak{b}\)生成，且等于\(\mathfrak{g}\)。（对该子代数应用 (a)，从而将问题化为\(\mathfrak{b}\)是\(\mathfrak{g}\)的理想的情形，并使用 § 1 的命题 4。）推出\(\mathfrak{g}\)的最少生成元数为\(\dim \mathfrak{g}/\mathcal{D}\mathfrak{g}\)。
\end{enumerate}

\begin{enumerate}
\def\labelenumi{\arabic{enumi}.}
\setcounter{enumi}{4}
\tightlist
\item
  \begin{enumerate}
  \def\labelenumii{(\alph{enumii})}
  \tightlist
  \item
    证明：若李代数 g 的每个子代数都是次不变的，则它是幂零的。（证明：若\(\dim g > 1\)，则每个\(x \in g\)都属于 g 的一个理想\(h \neq g\)，该理想具有与 g 相同的性质，因而由关于 g 维数的归纳假设可知它是幂零的；所以\(ad_{g} x\)是幂零的。）
  \end{enumerate}
\end{enumerate}

\begin{enumerate}
\def\labelenumi{(\alph{enumi})}
\setcounter{enumi}{1}
\tightlist
\item
  证明：若在李代数\(\mathfrak{g}\)中，每个不同于\(\mathfrak{g}\)的子代数都不同于其正规化子，则\(\mathfrak{g}\)是幂零的。（将其化归为 (a)。）
\end{enumerate}

\begin{enumerate}
\def\labelenumi{\arabic{enumi}.}
\setcounter{enumi}{5}
\item
  令 g 为幂零李代数，a 为 g 的交换理想。以下条件等价：(a) a 是 g 的极大交换理想；(b) a 是 g 的极大交换子代数；(c) a 等于其在 g 中的中心化子\(a'\)。（非 (a)\(\Rightarrow\)非 (b)\(\Rightarrow\)非 (c) 的蕴含显然。若\(a' \neq a\)，由命题 1 存在 g 的理想\(a''\)，满足\(a \subset a'' \subset a'\)且 dim\(a''/a = 1\)。于是\(a'' = a + Kx\)，从而\([a'', a''] \subset [a, a] + [x, a] = \{0\}\)，故 (a) 不成立。）
\item
  \begin{enumerate}
  \def\labelenumii{(\alph{enumii})}
  \tightlist
  \item
    令 V 为定义在\(K\)上的有限维向量空间，\(\mathfrak{g}\)为 V 上幂零自同态组成的李代数，\((\mathrm{V}_{r})_{0 \leq r \leq n}\)为 V 的递减向量子空间序列，满足\(V_{0} = V, V_{n} = \{0\}\)，且\(\mathfrak{g}(\mathrm{V}_{r}) \subset V_{r+1}\)在\(0 \leq r < n\)时成立。通过对 i 归纳证明\((\mathcal{D}^{i}\mathfrak{g})(\mathrm{V}_{r}) \subset V_{r+2^{i}}\)。若\(\dim V \geq 2^{i}\)，则 V 中被\(\mathcal{D}^{i}\mathfrak{g}\)湮灭的元素所成子空间的维数\(\geqslant 2^{i}\)。若\(\dim V \leqslant 2^{i}\)，则\(\mathcal{D}^{i}\mathfrak{g} = \{0\}\)。
  \end{enumerate}
\end{enumerate}

\begin{enumerate}
\def\labelenumi{(\alph{enumi})}
\setcounter{enumi}{1}
\item
  令\(\mathfrak{g}\)为幂零李代数，\(i\)为满足\(\geqslant 0\)的整数。若\(\dim \mathcal{D}^{i}\mathfrak{g} > 2^{i} + 1\)，则\(\mathcal{D}^{i}\mathfrak{g}\)的中心维数\(\geqslant 2^{i}\)。若\(\dim \mathcal{D}^{i}\mathfrak{g} \leqslant 2^{i} + 1\)，则\(\mathcal{D}^{i}\mathfrak{g}\)为交换的。（将 (a) 应用于\(\mathcal{D}^{i}\mathfrak{g}\)上的\(\operatorname{ad}_{\mathfrak{g}} x, x \in \mathfrak{g}\)的限制。）
\item
  令\(\mathfrak{g}\)为幂零李代数，\(i\)为满足\(\geqslant 0\)的整数。若\(\mathcal{D}^{i}\mathfrak{g}\)不交换，则\(\mathcal{D}^{i}\mathfrak{g}/\mathcal{D}^{i+1}\mathfrak{g}\)的维数\(\geqslant 2^{i} + 1\)。（将其转到商代数，化归为\(\dim \mathcal{D}^{i+1}\mathfrak{g} = 1\)的情形，并使用 (b)。）
\end{enumerate}

\begin{enumerate}
\def\labelenumi{\arabic{enumi}.}
\setcounter{enumi}{7}
\tightlist
\item
  令\(\mathfrak{g}\)为李代数，\(\mathfrak{m}\)为余维 1 的理想，\(x\)为\(\mathfrak{g}\)中不属于\(\mathfrak{m}\)的元素，且\(z \in \mathfrak{g}\)。
\end{enumerate}

\begin{enumerate}
\def\labelenumi{(\alph{enumi})}
\item
  证明线性映射 D 在 m 上为 0 并将 x 映到 z 时，若 z 属于 g 中 m 的中心化子 a，则 D 是导子。
\item
  令\(q\)为满足\(a \subset \mathcal{C}^q\mathfrak{g}\)的最大整数。证明若进一步有\(z \notin \mathcal{C}^{q+1}\mathfrak{g}\)，则\(D\)不是\(\mathfrak{g}\)的内导子。
\item
  由 (a)、(b) 和习题 7 (c) 推出：若 g 是维数为\textgreater1 的幂零李代数，则 g 的内导子空间在全体导子空间中的余维数为\(\geqslant2\)。
\end{enumerate}

\begin{enumerate}
\def\labelenumi{\arabic{enumi}.}
\setcounter{enumi}{8}
\tightlist
\item
  \begin{enumerate}
  \def\labelenumii{(\alph{enumii})}
  \tightlist
  \item
    验证下列乘法表定义了维数分别为 3 和 4 的两个幂零李代数\(\mathfrak{g}_3, \mathfrak{g}_4\)：
  \end{enumerate}
\end{enumerate}

\[
\mathfrak {g} _ {3}: \left[ x _ {1}, x _ {2} \right] = x _ {3}, \quad \left[ x _ {1}, x _ {3} \right] = \left[ x _ {2}, x _ {3} \right] = 0
\]

\[
\mathfrak {g} _ {4}: \left[ x _ {1}, x _ {2} \right] = x _ {3}, \quad \left[ x _ {1}, x _ {3} \right] = x _ {4}, \quad \left[ x _ {1}, x _ {4} \right] = \left[ x _ {2}, x _ {3} \right] = \left[ x _ {2}, x _ {4} \right] = \left[ x _ {3}, x _ {4} \right] = 0.
\]

\begin{enumerate}
\def\labelenumi{(\alph{enumi})}
\setcounter{enumi}{1}
\item
  证明\(\mathfrak{g}_3\)同构于\(\mathfrak{n}(3, \mathbf{K})\)，也同构于\(\mathfrak{sl}(2, \mathbf{K})\)；当\(\mathbf{K}\)的特征为 2 时上述结论成立。
\item
  令\(\mathfrak{g}_1\)为唯一的 1 维李代数。证明维数\(\leqslant 4\)的幂零李代数如下表所示：
\end{enumerate}

\[
\text { dimension } 1: \mathfrak {g} _ {1}; \text { dimension } 2: (\mathfrak {g} _ {1}) ^ {2};
\]

维数 3: \((\mathfrak{g}_1)^3, \mathfrak{g}_3\) ; 维数 4: \((\mathfrak{g}_1)^4, \mathfrak{g}_3 \times \mathfrak{g}_1, \mathfrak{g}_4\) .

（使用习题 7 (c)。若\(\mathfrak{g}\)是幂零的且维数为 3，并且\(\dim \mathcal{D}\mathfrak{g} = 1\)，注意到\(\mathcal{D}\mathfrak{g}\)包含于\(\mathfrak{g}\)的中心，从而\(\mathfrak{g} = \mathfrak{g}_3\)。若 dim \(\mathfrak{g} = 4\)，dim \(\mathcal{D}\mathfrak{g} = 1\)，注意括号运算在\(\mathfrak{g}\)上定义了\(\mathfrak{g} / \mathcal{D}\mathfrak{g}\)上的交错双线性形式；该形式必然退化，因此存在子代数\(\mathfrak{h}\)，它是\(\mathfrak{g}\)的子代数且包含于\(\mathfrak{g}\)的中心，满足 dim \(\mathfrak{h} = 1\)且\(\mathfrak{h}\cap \mathcal{D}\mathfrak{g} = \{0\}\)，从而\(\mathfrak{g} = \mathfrak{h}\times \mathfrak{g}_3\)。若 dim \(\mathfrak{g} = 4\)，dim \(\mathcal{D}\mathfrak{g} = 2\)，则\(\mathcal{D}\mathfrak{g}\)为交换的；将定理 1 应用于\(\mathcal{D}\mathfrak{g}\)上\(\operatorname{ad}_{\mathfrak{g}} x\)的限制\((x\in \mathfrak{g})\)，证明存在\(\mathfrak{h}\)的交换理想\(\mathfrak{g}\)，满足\(\mathfrak{h}\supseteq \mathcal{D}\mathfrak{g}\)且 dim \(\mathfrak{h} = 3\)；令\(x\in \mathfrak{g},\)，且\(x\notin \mathfrak{h}\)；选取\(\mathfrak{h}\)的一组基，使\(\operatorname{ad}_{\mathfrak{g}}x\)在\(\mathfrak{h}\)上的限制关于此基具有 Jordan 矩阵；于是\(\mathfrak{g} = \mathfrak{g}_4\)。

\begin{enumerate}
\def\labelenumi{\arabic{enumi}.}
\setcounter{enumi}{9}
\tightlist
\item
  令\(\mathcal{D}\)为\(\mathfrak{g}_4\)的导子李代数。证明\(\mathcal{D}\)的维数为 7、中心为零，并且内导子理想\(\mathcal{D}' \subset \mathcal{D}\)在\(\mathcal{D}\)中没有补子代数。
\end{enumerate}

¶ 11. 令 L 为交换环，A 为 L 上的左阿廷代数，γ 为从 A × A 到 L 的映射。在 A 上定义内部复合律\((a, b) \mapsto a * b = ab + \gamma(a, b)ba\)。对于 A 的每个子集 E，令\(\tilde{E}\)表示由 E 生成的 A 的子环（无单位元）。证明若 E 由幂零元组成且关于运算 * 稳定，则\(\tilde{E}\)是幂零的（即存在 n\textgreater{}0，使得\(\tilde{E}^{n} = \{0\}\)）。按以下步骤进行：

\begin{enumerate}
\def\labelenumi{\arabic{enumi}.}
\item
  首先设 A 是单环，因此同构于\(\mathcal{L}_{\mathrm{D}}(\mathrm{T})\)，其中 T 是域 D 上维数为\(m\)的左向量空间。对\(m\)作归纳。令\(\Phi\)为由满足以下条件的子集组成的集合：\(\mathrm{F} \subset \mathrm{E}\)关于\(*\)稳定，且\(\tilde{\mathrm{F}}\)是幂零的。证明\(\tilde{\Phi}\)具有极大元 M（注意，\(\tilde{\mathrm{F}}^m = \{0\}\)对所有\(\mathrm{F} \in \Phi\)都成立）。假设\(\tilde{\mathbf{M}} \neq \tilde{\mathbf{E}}\)。证明存在\(a \in \mathrm{E}\)，使得\(a \notin \tilde{\mathbf{M}}\)，且\(t * a \in \tilde{\mathbf{M}}\)对所有\(t \in \mathbf{M}\)都成立（注意，不存在无限序列\((a_n)\)，使得\(a_n \in \mathrm{E}\)、\(a_n \notin \mathbf{M}\)、\(a_n = t_{n-1} * a_{n-1}\)，且\(t_{n-1} \in \mathbf{M}\)）。令 S 为 T 的子空间，它是各个\(u(\mathrm{T})\)的直和，其中\(u \in \tilde{\mathbf{M}}\)；证明\(S \neq \{0\}\)且\(S \neq T\)，并且\(a(S) \subset S\)。令 N 为所有\(v \in E\)中满足\(v(S) \subset S\)者的集合。利用归纳假设，并把 N 的元素视为在 S 及 T/S 上的作用，证明\(N \in \Phi\)，从而矛盾。
\item
  一般情形下，利用 A 的 Jacobson 根基是幂零的这一事实。由此结果得到 Engel 定理的一个新证明。
\item
  令\(\mathfrak{g}\)为李代数，\(\mathcal{C}^{\infty}\mathfrak{g}\)为各个\(\mathcal{C}^p\mathfrak{g}\)的交。
\end{enumerate}

\begin{enumerate}
\def\labelenumi{(\alph{enumi})}
\item
 李代数\(\mathfrak{g} / \mathcal{C}^{\infty}\mathfrak{g}\)是幂零的。
\item
  证明存在幂零子代数\(\mathfrak{h}\)，它是\(\mathfrak{g}\)的子代数，满足\(\mathfrak{g} = \mathfrak{h} + \mathcal{C}^{\infty}\mathfrak{g}\)。（对\(\dim \mathfrak{g}\)作归纳。假设\(\mathfrak{g}\)不是幂零的。取\(x\in \mathfrak{g}\)，使得\(\mathbf{I} = \bigcap_{k}\left(\mathrm{(ad}x)^{k}(\mathfrak{g})\neq \{0\}\right)\)。令\(\mathfrak{n}\)为\((\mathrm{ad}x)^k\)的核的并。证明这是\(\mathfrak{g}\)的子代数，且\(\mathfrak{g}\)是\(\mathbf{I}\)与\(\mathfrak{n}\)的直和。由归纳假设，\(\mathfrak{n}\)是\(\mathcal{C}^{\infty}\mathfrak{n}\)与幂零子代数\(\mathfrak{h}\)之和。最后，\(\mathcal{C}^{\infty}\mathfrak{n}\subset \mathcal{C}^{\infty}\mathfrak{g}\)且\(\mathbf{I}\subset \mathcal{C}^{\infty}\mathfrak{g}\)。
\item
  验证下列乘法表定义了一个 5 维（可解）李代数 g：
\end{enumerate}

\[
\left[ x _ {1}, x _ {2} \right] = x _ {5}, \quad \left[ x _ {1}, x _ {3} \right] = x _ {3}, \quad \left[ x _ {1}, x _ {4} \right] = - x _ {4}, \quad \left[ x _ {3}, x _ {4} \right] = x _ {5}
\]

且其他\([x_i, x_j]\)均为零。证明对于此李代数

\[
\mathcal {C} ^ {\infty} \mathfrak {g} = \mathrm{K} x _ {3} + \mathrm{K} x _ {4} + \mathrm{K} x _ {5}
\]

但\(\mathcal{C}^{\infty}\mathfrak{g}\)在\(\mathfrak{g}\)中不存在补子代数。

¶ 13. 令 g 为李代数，\(\mathfrak{z}\)为 g 中\(C^{\infty}g\)的中心化子。若\(\mathfrak{z} \not\subset C^{\infty}g\)，则 g 的中心非零。（写成\(g = C^{\infty}g + t\)，其中 t 是 g 的幂零子代数（习题 12）。令\(g_{1} = \mathfrak{z} + t\)。令 b 为幂零子代数，满足\(g_{1} = C^{\infty}g_{1} + \mathfrak{b}\)。证明\(g = C^{\infty}g + \mathfrak{b}\)且\(C^{\infty}g_{1} \subset \mathfrak{z} \cap C^{\infty}g\)。令 y 为\(\mathfrak{z}\)中不属于\(C^{\infty}g\)的元素。写成\(y = x + x'\)，其中\(x \in \mathfrak{b}, x' \in C^{\infty}g_{1}\)；于是\(x \in \mathfrak{z}\)且\(x \neq 0\)，因此\(\mathfrak{z} \cap \mathfrak{b}\)是 b 的非零理想；利用习题 3 (a)，推出 g 中存在非零元素，它与\(C^{\infty}g\)及 b 可交换。）

¶ 14. 令 g 为李代数，b 为 g 的次不变子代数，且 b 在 g 中的中心化子\(\mathfrak{z}(b)\)为零。

\begin{enumerate}
\def\labelenumi{(\alph{enumi})}
\item
  中心化子\(\mathfrak{z}(\mathcal{C}^{\infty}\mathfrak{b})\)（即\(\mathcal{C}^{\infty}\mathfrak{b}\)在\(\mathfrak{g}\)中的中心化子）包含于\(\mathcal{C}^{\infty}\mathfrak{b}\)。（若\(\mathfrak{z}(\mathcal{C}^{\infty}\mathfrak{b})\not\subset\mathcal{C}^{\infty}\mathfrak{b},\mathfrak{z}(\mathcal{C}^{\infty}\mathfrak{b})\not\subset\mathfrak{b}\)，由习题 13 得；\(\mathcal{C}^{\infty}\mathfrak{b}\)是\(\mathfrak{g}\)的理想（§ 1，习题 14）；令\(\mathfrak{c}=\mathfrak{b}+\mathfrak{z}(\mathcal{C}^{\infty}\mathfrak{b})\)；\(\mathfrak{c}\)是子代数，\(\mathfrak{b}\)在\(\mathfrak{c}\)中是次不变的；\(\mathfrak{b}\)是\(\mathfrak{b}_{1}\subset\mathfrak{c}\)的理想，其中\(\mathfrak{b}\neq\mathfrak{b}_{1}\)；令\(y\in\mathfrak{b}_{1},y\notin\mathfrak{b};y=x+x'\)，其中\(x\in\mathfrak{z}(\mathcal{C}^{\infty}\mathfrak{b}),x'\in\mathfrak{b}\)；则\(x\in\mathfrak{z}(\mathcal{C}^{\infty}\mathfrak{b})\cap\mathfrak{b}_{1},x\notin\mathfrak{b};\mathfrak{b}_{1}=\mathrm{Kx}+\mathfrak{b}\)是子代数；\(\mathcal{C}^{\infty}\mathfrak{c}=\mathcal{C}^{\infty}\mathfrak{b}\)，因为\(\mathfrak{c}\cap\mathfrak{z}(\mathcal{C}^{\infty}\mathfrak{b})\subset\mathcal{C}^{\infty}\mathfrak{b}\)，从而\(\mathcal{C}^{k}\mathfrak{c}\subset\mathcal{C}^{k}\mathfrak{b}\)对所有\(k\)均成立；\(\mathfrak{z}(\mathcal{C}^{\infty}\mathfrak{b})\cap\mathfrak{b}\not\subset\mathcal{C}^{\infty}\mathfrak{b}\)，因此\(\mathfrak{b}\)的中心由习题 13 知非零；特别地，\(\mathfrak{z}(\mathfrak{b})\neq\{0\}\)，矛盾。）
\item
  由 (a) 推出，若\(\mathfrak{D}\)表示\(\mathcal{C}^{\infty}\mathfrak{b}\)的导子李代数，\(\mathfrak{a}\)表示\(\mathcal{C}^{\infty}\mathfrak{b}\)的中心，则\(\dim \mathfrak{g} \leqslant \dim \mathfrak{D} + \dim \mathfrak{a}\)（令\(\mathfrak{g}\)通过伴随表示作用于\(\mathcal{C}^{\infty}\mathfrak{b}\)。）
\end{enumerate}

\begin{enumerate}
\def\labelenumi{\arabic{enumi}.}
\setcounter{enumi}{14}
\tightlist
\item
  令\(l\)为中心为零的李代数，\(\mathfrak{D}_1\)为\(l\)的导子李代数，\(\mathfrak{D}_2\)为\(\mathfrak{D}_1, \ldots\)的导子李代数。于是\(l\)是\(\mathfrak{D}_1, \mathfrak{D}_1\)的理想，并且理想链继续到\(\mathfrak{D}_2\)，等等（§ 1，习题 15 (c)）。令\(\mathfrak{D}\)为\(\mathcal{C}^\infty l\)的导子李代数，\(a\)为\(\mathcal{C}^\infty l\)的中心。
\end{enumerate}

\begin{enumerate}
\def\labelenumi{(\alph{enumi})}
\item
  于是\(\dim \mathfrak{D}_1 \leqslant \dim \mathfrak{D} + \dim \mathfrak{a}\)。（使用上面的习题 14 和 § 1 的习题 15。）
\item
  由 (a) 推出，当 i 充分大时，\(D_{i}\)的所有导子都是内导子。
\end{enumerate}

\begin{enumerate}
\def\labelenumi{\arabic{enumi}.}
\setcounter{enumi}{15}
\item
  令\(\mathfrak{g}_1\)和\(\mathfrak{g}_2\)为李 K-代数，\(\mathfrak{n}_1\)和\(\mathfrak{n}_2\)分别为它们的最大幂零理想。证明\(\mathfrak{g}_1 \times \mathfrak{g}_2\)的最大幂零理想是\(\mathfrak{n}_1 \times \mathfrak{n}_2\)。
\item
  我们重新考察习题 9 (a) 中的李代数\(\mathfrak{g}_3\)，其基为\((x_1, x_2, x_3)\)。
\end{enumerate}

\begin{enumerate}
\def\labelenumi{(\alph{enumi})}
\tightlist
\item
  令 V 为向量空间 K{[}X{]}。令 D 为 V 上关于 X 的微分算子，M 为乘以 X 的算子。证明若 K 的特征为 0，则映射
\end{enumerate}

\[
\alpha x _ {1} + \beta x _ {2} + \gamma x _ {3} \mapsto \alpha D + \beta M + \gamma
\]

是 V 上的无限维简单表示\(\rho\)，其表示的李代数为\(\mathfrak{g}\)。

\begin{enumerate}
\def\labelenumi{(\alph{enumi})}
\setcounter{enumi}{1}
\tightlist
\item
  若 K 的特征为 p\textgreater{}0，则理想\((\mathbf{X}^{p})\)（即 K{[}X{]}的理想）在\(\rho(\mathfrak{g})\)下稳定。在 g 在\(\mathrm{K}[\mathrm{X}] / (\mathrm{X}^{p})\)上的商表示下，没有直线稳定。
\end{enumerate}

\begin{enumerate}
\def\labelenumi{\arabic{enumi}.}
\setcounter{enumi}{17}
\tightlist
\item
  令\(\mathfrak{g}\)为定义在\(\mathbf{K}\)上的 7 维向量空间，基为\((e_i)_{1\leqslant i\leqslant 7}\)。在\(\mathfrak{g}\)上按下列公式定义交错括号：
\end{enumerate}

\[
\left[ e _ {i}, e _ {j} \right] = \alpha_ {i j} e _ {i + j} \quad (1 \leqslant i <   j \leqslant 7, i + j \leqslant 7)\tag{1}
\]

且其他括号\([e_i, e_j]\)在\(i < j\)时均为零。

\begin{enumerate}
\def\labelenumi{(\alph{enumi})}
\tightlist
\item
  要使 Jacobi 恒等式成立，充要条件是：
\end{enumerate}

\begin{enumerate}
\def\labelenumi{(\arabic{enumi})}
\setcounter{enumi}{1}
\tightlist
\item
\end{enumerate}

\[
- \alpha_ {2 3} \alpha_ {1 5} + \alpha_ {1 3} \alpha_ {2 4} = 0\tag{2}
\]

\[
\alpha_ {1 2} \alpha_ {3 4} - \alpha_ {2 4} \alpha_ {1 6} + \alpha_ {1 4} \alpha_ {2 5} = 0.\tag{3}
\]

\begin{enumerate}
\def\labelenumi{(\alph{enumi})}
\setcounter{enumi}{1}
\item
  以下设所有\(\alpha_{ij}\)均\(\neq 0\)。证明理想\(\mathcal{C}^2\mathfrak{g}\)、\(\mathcal{C}^3\mathfrak{g}\)、\(\mathcal{C}^4\mathfrak{g}\)、\(\mathcal{C}^5\mathfrak{g}\)、\(\mathcal{C}^6\mathfrak{g}\)分别具有以下基：\((e_3, e_4, e_5, e_6, e_7)\)、\((e_4, e_5, e_6, e_7)\)、\((e_5, e_6, e_7)\)、\((e_6, e_7)\)、\((e_7)\)。证明\(\mathfrak{h}\)的中心化子\(\mathcal{C}^5\mathfrak{g}\)具有基\((e_2, e_3, e_4, e_5, e_6, e_7)\)。
\item
  令\((e_i')_{1 \leqslant i \leqslant 7}\)为\(\mathfrak{g}\)的另一组基，满足\(\mathcal{C}^6\mathfrak{g}\)具有基\((e_7')\)，\(\mathcal{C}^5\mathfrak{g}\)具有基\((e_6', e_7')\)，\ldots\(\mathfrak{h}\)具有基\((e_2', e_3', e_4', e_5', e_6', e_7')\)。证明
\end{enumerate}

\[
\left[ e _ {i} ^ {\prime}, e _ {j} ^ {\prime} \right] = \alpha_ {i j} ^ {\prime} e _ {i + j} ^ {\prime} + \beta e _ {i + j + 1} ^ {\prime} + \gamma e _ {i + j + 2} ^ {\prime} + \dots + \lambda e _ {7} ^ {\prime},
\]

其中

\[
\alpha_ {1 4} ^ {\prime} \alpha_ {2 5} ^ {\prime} \alpha_ {1 6} ^ {\prime - 1} \alpha_ {2 4} ^ {\prime - 1} = \alpha_ {1 4} \alpha_ {2 5} \alpha_ {1 6} ^ {- 1} \alpha_ {2 4} ^ {- 1}.
\]

\begin{enumerate}
\def\labelenumi{(\alph{enumi})}
\setcounter{enumi}{3}
\tightlist
\item
  由此推出：若 K 是无限域，则存在无穷多个互不同构的 7 维 K 上幂零李代数；并且存在不能通过将基域从 R 扩张到 C 而由实幂零李代数得到的复幂零李代数。
\end{enumerate}

\begin{enumerate}
\def\labelenumi{\arabic{enumi}.}
\setcounter{enumi}{18}
\tightlist
\item
  \begin{enumerate}
  \def\labelenumii{(\alph{enumii})}
  \tightlist
  \item
    令\(\mathfrak{g}\)为李代数。令\(\mathfrak{S}\)为\(\mathfrak{g}\)的导子李代数。\(\mathfrak{g}^{[k]}\)的特征理想（在\(\mathfrak{g}\)中）按如下归纳方式定义：\(\mathfrak{g}^{[0]} = \mathfrak{g}\)，且\(\mathfrak{g}^{[k+1]}\)是由\(\mathfrak{g}\)中的 Dx 生成的子空间（\(\mathrm{D} \in \mathfrak{S}, x \in \mathfrak{g}^{[k]}\)）。以下条件等价：(1)\(\mathfrak{g}\)的每个导子都是幂零的；(2)当\(\mathfrak{g}^{[k]} = \{0\}\)在\(k\)充分大时成立；(3)\(\mathfrak{g}\)的全纯代数是幂零的。若这些条件成立，则\(\mathfrak{g}\)称为特征幂零。这样的代数是幂零的。
  \end{enumerate}
\end{enumerate}

\begin{enumerate}
\def\labelenumi{(\alph{enumi})}
\setcounter{enumi}{1}
\tightlist
\item
  验证下述乘法表定义了一个 8 维李代数：
\end{enumerate}

\[
\begin{array}{r} \left[ x _ {1}, x _ {2} \right] = x _ {3}, \quad \left[ x _ {1}, x _ {3} \right] = x _ {4}, \quad \left[ x _ {1}, x _ {4} \right] = x _ {5}, \quad \left[ x _ {1}, x _ {5} \right] = x _ {6} \\ \left[ x _ {1}, x _ {6} \right] = x _ {8}, \quad \left[ x _ {1}, x _ {7} \right] = x _ {6}, \quad \left[ x _ {2}, x _ {3} \right] = x _ {5}, \quad \left[ x _ {2}, x _ {4} \right] = x _ {6} \\ \left[ x _ {2}, x _ {5} \right] = x _ {7}, \quad \left[ x _ {2}, x _ {6} \right] = 2 x _ {8}, \quad \left[ x _ {3}, x _ {4} \right] = - x _ {7} + x _ {8}, \quad \left[ x _ {3}, x _ {5} \right] = - x _ {8} \end{array}
\]

且\([x_i, x_j] = 0\)在\(i + j > 8\)时成立；证明

\[
\begin{array}{c c} \mathcal {C} ^ {2} \mathfrak {g} = \sum_ {i = 3} ^ {8} \mathrm{K} x _ {i}, & \mathcal {C} ^ {3} \mathfrak {g} = \sum_ {i = 4} ^ {8} \mathrm{K} x _ {i}, \\ & \mathcal {C} ^ {6} \mathfrak {g} = \mathrm{K} x _ {8}, \quad \mathcal {C} ^ {7} \mathfrak {g} = \{0 \}, \end{array} \quad \begin{array}{c c} \mathcal {C} ^ {4} \mathfrak {g} = \sum_ {i = 5} ^ {8} \mathrm{K} x _ {i}, & \mathcal {C} ^ {5} \mathfrak {g} = \sum_ {i = 6} ^ {8} \mathrm{K} x _ {i}, \\ & [ \mathcal {C} ^ {2} \mathfrak {g}, \mathcal {C} ^ {2} \mathfrak {g} ] = \mathrm{K} x _ {7} + \mathrm{K} x _ {8}, \end{array}
\]

并且\(\mathcal{C}^2\mathfrak{g}\)到\(\mathcal{C}^4\mathfrak{g}\)的输运子空间是\(\sum_{i=2}^{\infty}\mathrm{Kx}_i\)。推出每个\(\mathfrak{g}\)的导子\(\mathrm{D}x_i = \sum_{j=1}^{8}u_{ij}x_j\)均由公式定义。再证明：当 K 的特征不为 2 时，\(u_{ii} = 0\)对所有\(i\)均成立，从而\(\mathfrak{g}\)是特征幂零的。

\begin{enumerate}
\def\labelenumi{(\alph{enumi})}
\setcounter{enumi}{2}
\tightlist
\item
  若 K 的特征为\(\neq2\)，证明 (b) 中的李代数 g 不是任何李代数的导代数。（若 g = Dh，先证明 h 是幂零的。注意\(\dim g/Dg = 2\)；结合习题 7 (c) 得到矛盾。）
\end{enumerate}

¶ 20. 令 g 为李代数，U 为其包络代数。

\begin{enumerate}
\def\labelenumi{(\alph{enumi})}
\item
  令\(\mathfrak{g}'\)为\(\mathfrak{g}\)的理想，\(\mathrm{U}' \subset \mathrm{U}\)为其包络代数，\(x \in \mathfrak{g}\)以及\(a_1, \ldots, a_n\)为 U 中的元素。假设\(a_i = x\)对指标\(j_1, \ldots, j_n\)成立，并且\(a_i \in \mathrm{U}'\)对其他指标成立（令\(k_1, \ldots, k_q\)为这些指标，且满足\(k_1 < k_2 < \cdots < k_q\)）。那么\(a_1 a_2 \ldots a_n - x^p a_{k_1}a_{k_2} \ldots a_{k_q}\)是形如\(x^p b\)的项之和，其中\(b \in \mathrm{U}'\)且\(p' < p\)。（对\(p\)作归纳。）
\item
  此后设 g 为幂零李代数，K 的特征为 0。令 g' 为 g 中余维 1 的理想，U' 为其包络代数，x 为 g 中不属于 g' 的元素，U₀（分别为 U₀'）为 U（分别为 U'）中被 g 的一组导子 b 湮灭的子代数（这些导子延拓为 U 的导子，并将 g 映入 g'）。假设 U₀ 包含于 U 的中心且 U₀ ≠ U'。证明存在 a₁ ∈ U₀'、a₂ ∈ U'，使得 a₁ ≠ 0 且 a = xa₁ + a₂ ∈ U₀。（令 xᵐbₘ + xᵐ⁻¹bₘ₋₁ + ⋯ + b₀ ∈ U₀，其中 bₘ、\ldots、b₀ 属于 U'，且 m\textgreater{}0，bₘ ≠ 0。利用⟨a⟩证明：对 g 的每个导子 D ∈ b，都有 Dbₘ = 0，m(Dx)bₘ + Dbₘ₋₁ = 0，从而 D(mxbₘ + bₘ₋₁) = 0。）证明 U₀ 包含于其分式域中的代数 K{[}a, a₁⁻¹, U₀'{]}，该代数由 a、a₁⁻¹、U₀' 生成（可由 § 2 定理 1 的推论 7 构造）。推出 U₀ 的分式域是由 a 和 U₀' 生成的域。证明 a 在 K(U₀') 上超越。
\item
 令\(\{0\} = \mathfrak{g}_0 \subset \mathfrak{g}_1 \subset \cdots \subset \mathfrak{g}_n = \mathfrak{g}\)为\(\mathfrak{g}\)的维数分别为\(0, 1, \ldots, n\)的理想序列。令\(x_j\)为\(\mathfrak{g}_j\)中不属于\(\mathfrak{g}_{j-1}\)的元素。令\(j_1 < j_2 < \cdots < j_n\)为满足以下条件的指标\(j\)：存在\(\mathrm{U}^j\)（\(\mathfrak{g}_j\)的包络代数）中的\(\mathrm{U}_0\)（\(\mathrm{U}\)的中心），它不属于\(\mathrm{U}^{j-1}\)。根据 (b)，存在\(a_{j_1} \in \mathrm{U}_0 \cap \mathrm{U}^{j_1}\)和\(a_{j_2} \in \mathrm{U}^{j_2-1}\)，满足\(a_{j_1} \neq 0\)且\(a_j = x_j a_{j_1} + a_{j_2} \in \mathrm{U}_0 \cap \mathrm{U}^{j}\)。对\(n\)作归纳证明：
\end{enumerate}

\[
\mathrm{U} _ {0} \subset \mathrm{K} [ a _ {j _ {1}}, \dots , a _ {j _ {q}}, a _ {j _ {1}} ^ {- 1}, \dots , a _ {j _ {q}} ^ {- 1} ]
\]

并且\(U_{0}\)的分式域由代数独立元\(a_{j_{1}},\ldots,a_{j_{q}}\)生成。特别地，U 中心的分式域是 K 的纯超越扩张。

\begin{enumerate}
\def\labelenumi{\arabic{enumi}.}
\setcounter{enumi}{20}
\tightlist
\item
  令 g 为李代数，σ 为 g 的自同构。
\end{enumerate}

\begin{enumerate}
\def\labelenumi{(\alph{enumi})}
\tightlist
\item
  首先设\(\mathbf{K}\)代数闭。对任意\(\lambda \in \mathbf{K}\)，令\(\mathrm{L}_{\lambda}\)为由\(x \in \mathfrak{g}\)中被\(\sigma - \lambda \mathrm{I}\)的幂所湮灭的元素组成的集合；\(\mathfrak{g}\)是各个\(\mathrm{L}_{\lambda}\)的直和。证明\([\mathrm{L}_{\lambda}, \mathrm{L}_{\mu}] \subset \mathrm{L}_{\lambda \mu}\)。（注意
\end{enumerate}

\[
(\sigma - \lambda \mu \mathrm{I}) ([ x, y ]) = [ (\sigma - \lambda \mathrm{I}) x, \sigma y ] + [ \lambda x, (\sigma - \mu \mathrm{I}) y ].)
\]

\begin{enumerate}
\def\labelenumi{(\alph{enumi})}
\setcounter{enumi}{1}
\item
  对任意域 K，设\(\sigma\)的任一特征值（在 K 的代数闭扩张中）都不是单位根。证明 g 是幂零的。（将其化归为 K 代数闭的情形。若\(\lambda_{1},\ldots,\lambda_{m}\)是\(\sigma\)的不同特征值，则\(\lambda_{i}\lambda_{j},\lambda_{i}\lambda_{j}^{2},\lambda_{i}\lambda_{j}^{3},\ldots\)不全是特征值，因此\(ad_{g}x\)对\(x\in L_{\lambda_{i}}\)是幂零的。使用习题 11，将\(ad_{g}x\)应用于其中 x 遍历\(L_{\lambda_{i}}\)的集合。）
\item
  对任意域 K，设\(\sigma^{q}=I\)，其中 q 为素数，且\(\sigma\)没有特征值等于 1。证明 g 是幂零的。（采用 (b) 中的方法，注意：若\(g\neq\{0\}\)，则 q 不等于 K 的特征，并且对\(\lambda_{i},\lambda_{j}\)中的每一对有序特征值\(\sigma\)，存在整数 k 使\(\lambda_{i}\lambda_{j}^{k}=1\)。）
\end{enumerate}

\begin{enumerate}
\def\labelenumi{\arabic{enumi}.}
\setcounter{enumi}{21}
\tightlist
\item
  \begin{enumerate}
  \def\labelenumii{(\alph{enumii})}
  \tightlist
  \item
    令\(u, v\)为有限维向量空间\(\mathbf{E}\)上的两个自同态，其中该空间定义在代数闭域\(\mathbf{K}\)上。对\(\lambda\)（\(u\)的特征值），令\(\mathrm{E}_{\lambda}\)为\(\mathbf{E}\)中被某个\(u - \lambda \mathrm{I}\)的幂所湮灭的向量组成的子空间。证明若\((\operatorname{ad} u)^n v = 0\)，则\(\mathrm{E}_{\lambda}\)在\(v\)下稳定。
  \end{enumerate}
\end{enumerate}

\begin{enumerate}
\def\labelenumi{(\alph{enumi})}
\setcounter{enumi}{1}
\item
  由 (a) 推出，若 g 是 E 上自同态组成的幂零李代数，则 E 是子空间\(F_{j}\)的直和（\(1 \leqslant j \leqslant m\)），这些子空间在 g 下稳定，并且在每个\(F_{j}\)上，每个\(u \in g\)的限制都可写成\(\lambda_{j}(u)\mathbf{I} + u_{j}\)，其中\(\lambda_{j}(u) \in \mathbf{K}\)，且\(u_{j}\)为幂零的。
\item
  若\(\mathbf{K}\)的特征为 2，\(\mathrm{E} = \mathrm{K}^2\)，且\(\mathfrak{g}\)为幂零代数\(\mathfrak{sl}(2,\mathbf{K})\)，证明\(m = 1\)，并且可能有\(\lambda (u + v)\neq \lambda (u) + \lambda (v)\)，其中\(u,v\)是\(\mathfrak{g}\)中的两个元素。（关于本题的后续内容，参见 § 5，习题 12。）
\end{enumerate}

\begin{enumerate}
\def\labelenumi{\arabic{enumi}.}
\setcounter{enumi}{22}
\tightlist
\item
  令\(\mathfrak{g}\)为李 \(p\)-代数，定义在完美域\(\mathbf{K}\)上，且该域的特征为\(p > 0\)。\(\mathfrak{g}\)称为\(p\)-幂幺的，如果对所有\(x \in \mathfrak{g}\)，存在\(m\)使\(x^{p^m} = 0\)。
\end{enumerate}

\begin{enumerate}
\def\labelenumi{(\alph{enumi})}
\item
  证明每个\(p\)-幂幺李 \(p\)-代数都是幂零的。
\item
  设\(\mathfrak{g}\)为幂零的，并令\(\mathfrak{h}\)表示\(p\)-核，取自\(\mathfrak{g}\)的中心（§ 1，习题 23 (a)）。证明\(\mathfrak{g} / \mathfrak{h}\)是\(p\)-幂幺的。
\item
  令\(\mathfrak{g}\)为幂零李 \(p\)-代数，其三个基元素为\(e_1, e_2, e_3\)，满足\([e_1, e_2] = [e_1, e_3] = 0\)、\([e_2, e_3] = e_1\)、\(e_1^p = e_1\)和\(e_2^p = e_3^p = 0\)。则\(\mathfrak{h} = \mathrm{Ke}_1\)，但\(\mathfrak{g}\)不是\(\mathfrak{h}\)与一个\(p\)-幂幺\(p\)-子代数的直和。
\item
  若 g 是 p-幂幺的，证明在 g 的受限包络代数中，由 g 生成的双边理想是幂零的。（对 g 的维数作归纳。）
\end{enumerate}

\begin{enumerate}
\def\labelenumi{\arabic{enumi}.}
\setcounter{enumi}{23}
\item
  假设 K 的特征为 2。证明习题 9 中的幂零李代数\(\mathfrak{g}_4\)不存在 2-映射。
\item
  令\(\mathfrak{g}\)为李 \(p\)-代数。证明\(\mathfrak{g}\)的最大幂零理想是\(p\)-理想（参见 § 1，习题 22）。
\item
  令\(\mathbf{G}\)为一个群，令\((\mathbf{H}_n)_{n\geqslant 1}\)为\(\mathbf{G}\)的子群递减序列；设\(\mathbf{H}_1 = \mathbf{G}\)，并且若记\((x,y) = xyx^{-1}y^{-1}\)，则关系\(x\in \mathrm{H}_i,y\in \mathrm{H}_j\)蕴含\((x,y)\in \mathrm{H}_{i + j}\)。
\end{enumerate}

\begin{enumerate}
\def\labelenumi{(\alph{enumi})}
\item
  令\(G_{i} = H_{i}/H_{i+1}\)；证明\(G_{i}\)为交换的，并且映射\(x, y \mapsto (x, y)\)在取商后定义了从\(G_{i} \times G_{j}\)到\(G_{i+j}\)的 Z-双线性映射。
\item
  令\(\mathrm{gr}(\mathbf{G}) = \sum_{i=1}^{\infty} \mathbf{G}_i\)，并将 (a) 中定义的映射\(\mathbf{G}_i \times \mathbf{G}_j \to \mathbf{G}_{i+j}\)按线性延拓为从\(\mathrm{gr}(\mathbf{G}) \times \mathrm{gr}(\mathbf{G})\)到\(\mathrm{gr}(\mathbf{G})\)的 Z-双线性映射。证明\(\mathrm{gr}(\mathbf{G})\)关于此映射是李 Z-代数（验证 Jacobi 恒等式时，使用下列公式：
\end{enumerate}

\[
((x, y), z ^ {y}) \cdot ((y, z), x ^ {z}) \cdot ((z, x), y ^ {x}) = e \quad x, y, z \text {   in   } G,
\]

其中\(x^{y}\)表示\(yxy^{-1}\)，\(e\)表示 G 的单位元）。

\begin{enumerate}
\def\labelenumi{(\alph{enumi})}
\setcounter{enumi}{2}
\tightlist
\item
  若存在\(n\)使得\(\mathrm{H}_n = \{e\}\)，证明\(\operatorname{gr}(G)\)是幂零李 \(\mathbf{Z}\)-代数。
\end{enumerate}

\begin{enumerate}
\def\labelenumi{\arabic{enumi}.}
\setcounter{enumi}{26}
\tightlist
\item
   令 A 为带单位元 1 的结合代数，并令\(\mathrm{A}_0 = \mathrm{A} \supset \mathrm{A}_1 \supset \dots \supset \mathrm{A}_n \supset \dots\)为 A 的双边理想递减序列，满足\(\mathrm{A}_i \cdot \mathrm{A}_j \subset \mathrm{A}_{i+j}\)。令 G 为单位元\(e\)的群，并令\(f: \mathrm{G} \to \mathrm{A}\)为映射，满足\(f(e) = 1\)、\(f(xy) = f(x) \cdot f(y)\)，且\(1 - f(x) \in \mathrm{A}_1\)对所有\(x \in \mathrm{G}\)成立。令\(\mathrm{H}_n\)表示所有\(x \in \mathrm{G}\)中满足\(1 - f(x) \in \mathrm{A}_n\)者的集合。证明\(\mathrm{H}_n\)满足习题 26 的条件。证明映射\(x \mapsto f(x) - 1\)在取商后定义了李代数\(\mathrm{gr}(\mathrm{G})\)到与分次环\(\mathrm{gr}(\mathrm{A}) = \sum \mathrm{A}_n / \mathrm{A}_{n+1}\)相关的李代数中的单射同态（参见《交换代数》第三章）。
\end{enumerate}

% semantic-fragments: legacy-input-seam
\relax{}
